\documentclass[11pt,a4paper]{article}

\usepackage[utf8]{inputenc}
\usepackage[T1]{fontenc}
\usepackage{lmodern}
\usepackage[a4paper,top=2.4cm,bottom=2.4cm,left=2.6cm,right=2.6cm]{geometry}
\usepackage{parskip}
\usepackage{enumitem}
\usepackage{xcolor}
\usepackage{fancyhdr}
\usepackage{titlesec}

\definecolor{polito}{RGB}{0,60,113}
\definecolor{rulegray}{RGB}{200,200,200}

\titleformat{\section}{\normalfont\large\bfseries\color{polito}}{\thesection.}{0.5em}{}
\titlespacing*{\section}{0pt}{1.6\baselineskip}{0.5\baselineskip}

\pagestyle{fancy}
\fancyhf{}
\renewcommand{\headrulewidth}{0.4pt}
\fancyhead[L]{\footnotesize Scheda di sintesi --- Dylan Magliano (s337915)}
\fancyfoot[C]{\footnotesize \thepage}

\newcommand{\campo}[2]{%
  \par\noindent\parbox[t]{0.30\textwidth}{\raggedright\bfseries #1}%
  \parbox[t]{0.68\textwidth}{\raggedright #2}\par\medskip}

\newcommand{\camposdg}[2]{%
  \par\noindent\parbox[t]{0.22\textwidth}{\raggedright\bfseries #1}%
  \parbox[t]{0.76\textwidth}{\raggedright #2}\par\medskip}

\begin{document}

\begin{center}
  {\Large\bfseries\color{polito} POLITECNICO DI TORINO}\\[0.35em]
  {\normalsize Corso di Laurea Magistrale in Ingegneria Informatica}\\
  {\small\itshape Orientamento Artificial Intelligence and Data Analytics}\\[1.2em]
  {\large\bfseries Scheda di sintesi della tesi di laurea magistrale}\\[0.4em]
  {\small Anno Accademico 2025--2026}
\end{center}

\vspace{0.6em}
\noindent\textcolor{rulegray}{\rule{\textwidth}{0.8pt}}
\vspace{0.4em}

\campo{Candidato}{Dylan Magliano --- Matricola s337915}
\campo{Relatori}{prof. Flavio Giobergia\\ dott. Sergio Caprioli\\ dott. Emanuele Cagliero}
\campo{Azienda}{Intesa Sanpaolo}

\noindent\textcolor{rulegray}{\rule{\textwidth}{0.8pt}}

\section{Titolo della tesi (italiano)}

\textbf{Dai pixel ai portafogli finanziari: generare scenari di mercato
realistici tramite variational autoencoder}

\section{Thesis title (English)}

\textbf{From pixels to financial portfolios: generating realistic market
scenarios with variational autoencoders}

\section{Abstract in italiano}

{\small\itshape Lunghezza: \input{count_it.tex}\unskip{} caratteri spazi
inclusi (limite 3550).\par}
\vspace{0.4em}

La stima delle matrici di correlazione fra asset finanziari a partire dalle sole serie storiche presenta due limiti strutturali: il rumore statistico che degrada le matrici empiriche quando il numero di titoli è dello stesso ordine di grandezza della lunghezza della serie, e il fatto che il passato offra una sola traiettoria fra tutte quelle possibili. Questa tesi, svolta in collaborazione con Intesa Sanpaolo, affronta entrambi i limiti testando l'efficacia dei Variational Autoencoder (VAE) come modello generativo stocastico per la simulazione di scenari di mercato realistici.

Il lavoro parte dai rendimenti giornalieri di 362 titoli dell'indice S\&P 500 quotati con continuità nel ventennio 2000-2020. Da essi vengono estratte, con finestra mobile di 724 giorni e passo di 10, 461 matrici di correlazione empiriche, ridotte a 412 da un filtro di validità basato sulla proprietà di Perron-Frobenius. Poiché una matrice prodotta da un modello generativo non è automaticamente semidefinita positiva --- vincolo imprescindibile per una matrice di correlazione valida --- l'input dei modelli non è la matrice grezza, bensì il fattore triangolare inferiore della sua decomposizione di Cholesky: qualunque vettore prodotto dal decoder si ricompone così in una matrice semidefinita positiva per costruzione.

La metodologia segue un percorso incrementale su quattro architetture confrontate a parità di dimensione latente: Principal Component Analysis come benchmark lineare, Autoencoder Lineare, Autoencoder Profondo e Variational Autoencoder. La valutazione combina metriche di errore spaziale (MSE, MAE, norma di Frobenius) e metriche topologiche costruite sul Minimum Spanning Tree. Sul piano della pura ricostruzione storica l'Autoencoder Profondo domina tutti i modelli a ogni dimensione latente; il VAE, addestrato con una loss di ricostruzione basata sulla Divergenza di Jeffreys fra le distribuzioni associate alle matrici, si colloca in posizione intermedia: a venti dimensioni latenti supera sia la PCA sia l'Autoencoder Lineare, pur restando dietro all'Autoencoder deterministico. La regolarizzazione variazionale non compromette dunque in modo sostanziale la fedeltà di ricostruzione.

È sul piano generativo che il VAE mostra il proprio valore distintivo: la continuità e la densità dello spazio latente indotte dalla Divergenza di Kullback-Leibler sono precluse a un autoencoder deterministico, il cui spazio latente resta frammentario. La procedura adottata non è un semplice campionamento dal prior, bensì un block-bootstrap storico degli incrementi latenti giornalieri, che ancora ciascuno scenario allo stato di mercato osservato oggi e produce una distribuzione di evoluzioni plausibili della struttura di correlazione a un orizzonte di un mese, valutabile in ottica di rischiosità.

I mille scenari generati sono stati validati rispetto ai cinque fatti stilizzati delle matrici di correlazione finanziarie: distribuzione dei coefficienti spostata verso valori positivi, spettro degli autovalori conforme alla legge di Marchenko-Pastur, proprietà di Perron-Frobenius, struttura gerarchica settoriale e topologia scale-free del Minimum Spanning Tree. Tutti e mille gli scenari soddisfano simultaneamente i cinque criteri, con similarità coseno superiore a 0.98 rispetto alla matrice odierna sui modi principali. L'architettura proposta si conferma quindi uno strumento autonomo per la generazione di ecosistemi finanziari sintetici e matematicamente coerenti.


\section{Abstract in English}

{\small\itshape Length: \input{count_en.tex}\unskip{} characters including
spaces (limit 3550).\par}
\vspace{0.4em}

Estimating correlation matrices between financial assets from historical series alone suffers from two structural limitations: the statistical noise that degrades empirical matrices when the number of assets is of the same order of magnitude as the length of the series, and the fact that the past offers a single trajectory among all those that were possible. This thesis, carried out in collaboration with Intesa Sanpaolo, addresses both limitations by testing the effectiveness of Variational Autoencoders (VAE) as stochastic generative models for the simulation of realistic market scenarios.

The work starts from the daily returns of 362 S\&P 500 constituents continuously quoted over the twenty-year period 2000-2020. From them, using a rolling window of 724 days with a stride of 10, 461 empirical correlation matrices are extracted, reduced to 412 by a validity filter based on the Perron-Frobenius property. Since a matrix produced by a generative model is not automatically positive semidefinite --- a necessary condition for any valid correlation matrix --- the input to the models is not the raw matrix but the lower triangular factor of its Cholesky decomposition: any vector produced by the decoder therefore recomposes into a positive semidefinite matrix by construction.

The methodology follows an incremental path across four architectures compared at equal latent dimension: Principal Component Analysis as a linear benchmark, Linear Autoencoder, Deep Autoencoder and Variational Autoencoder. The evaluation combines spatial error metrics (MSE, MAE, Frobenius norm) with topological metrics built on the Minimum Spanning Tree. On pure historical reconstruction the Deep Autoencoder achieves the best results at every latent dimension considered; the VAE, trained with a reconstruction loss based on the Jeffreys Divergence between the distributions associated with the matrices, sits in an intermediate position: at twenty latent dimensions it outperforms both PCA and the Linear Autoencoder, while remaining behind the deterministic Autoencoder. Variational regularisation therefore does not substantially compromise reconstruction fidelity.

It is on the generative side that the VAE shows its distinctive value: the continuity and density of the latent space induced by the Kullback-Leibler Divergence are not guaranteed in a deterministic autoencoder, whose latent space remains fragmented. The procedure adopted is not a simple sampling from the prior, but a historical block bootstrap of the daily latent increments, which anchors each scenario to the market state observed today and produces a distribution of plausible one-month evolutions of the correlation structure, to be assessed from a risk perspective.

The one thousand generated scenarios were validated against the five stylized facts of financial correlation matrices: a coefficient distribution shifted towards positive values, an eigenvalue spectrum consistent with the Marchenko-Pastur law, the Perron-Frobenius property, hierarchical sector structure and the scale-free topology of the Minimum Spanning Tree. All one thousand scenarios satisfy the five criteria simultaneously, with cosine similarity above 0.98 with respect to today's matrix on the principal modes. The proposed architecture thus emerges as a self-contained and structurally coherent tool for generating realistic synthetic financial ecosystems.


\section{Keywords (English)}

\begin{itemize}[leftmargin=1.3em,itemsep=0.15em,topsep=0.3em]
  \item Variational autoencoder
  \item Deep generative models
  \item Synthetic scenario generation
  \item Financial correlation matrices
  \item Cholesky decomposition
  \item Positive semidefinite matrices
  \item Stylized facts
  \item Random matrix theory
  \item Marchenko--Pastur distribution
  \item Minimum spanning tree
  \item Latent space block bootstrap
  \item Dimensionality reduction
  \item Representation learning
  \item Quantitative finance
  \item Market risk assessment
  \item S\&P 500
\end{itemize}

\section{Sustainable Development Goals (SDG)}

\camposdg{SDG prevalente}{\textbf{09. INDUSTRY, INNOVATION AND INFRASTRUCTURE}}

\noindent\small Il contributo centrale della tesi è metodologico e
tecnologico: progetta e valuta un'architettura di deep learning generativo che
estende alle matrici di correlazione finanziarie strumenti nati per il dominio
delle immagini. Rientra quindi nel target 9.5, che chiede di potenziare la
ricerca scientifica e di aggiornare le capacità tecnologiche dei settori
industriali, incoraggiando l'innovazione.\normalsize

\vspace{1.2em}

\camposdg{SDG secondario}{\textbf{08. DECENT WORK AND ECONOMIC GROWTH}}

\noindent\small Il lavoro nasce in collaborazione con Intesa Sanpaolo e
fornisce agli istituti finanziari uno strumento per generare scenari di
mercato realistici con cui stressare i propri modelli. Si collega al target
8.10 (rafforzare la capacità delle istituzioni finanziarie): una migliore
rappresentazione della struttura di rischio sostiene la stabilità del sistema
creditizio e, per suo tramite, la crescita economica.\normalsize

\vspace{1.2em}

\camposdg{SDG terziario}{\textbf{10. REDUCED INEQUALITIES}}

\noindent\small Il target 10.5 chiede di migliorare la regolamentazione e il
monitoraggio dei mercati e delle istituzioni finanziarie globali. Metodi di
generazione di scenari validati su proprietà strutturali verificabili (i
cinque fatti stilizzati) rendono più trasparente e controllabile la
rappresentazione del rischio, presupposto tecnico di una vigilanza efficace e
di mercati più stabili.\normalsize

\end{document}
